% Scope: elementary diagnostic results, not a new general theory of forgetting.
% Repository derivations: doc/auto/theory-propositions.md; inherited drafts remain unchanged.
\section{What collapse diagnostics can identify}
\label{sec:theory}

An ordering intervention changes the sequence of state transitions. Whether it
changes memory depends on the cue-bearing channel, the decoder, and the noise
model. The following elementary results separate these dependencies. They
organize established linear-memory and attractor arguments
\citep{ganguli2008-memory-traces,seung1996,pals2024-phase-locked-memory}; no
novelty is claimed for the individual results. Proofs and additional
counterexamples appear in \cref{app:proofs}.

\subsection{Ordering nulls concern the cue channel}

Let $E(u)\in\R^n$ be the state after writing cue $u$, and let
$h=(x_0,\ldots,x_{H-1})$ be a distractor body independent of $u$.
All comparisons use the same encoder and body length $H$.

\begin{proposition}[Affine cue-channel ordering nulls]
\label{prop:affine}\label{prop:commuting}
For $s_{t+1}=As_t+c_{x_t}$ with $s_0=E(u)$,
\begin{equation}
 s_H(u,h)=A^H E(u)+k_h,\qquad
 k_h=\sum_{j=0}^{H-1}A^{H-1-j}c_{x_j}.
 \label{eq:affine-channel}
\end{equation}
Consequently cue differences, cue derivatives (when defined), and the
body-conditioned cue covariance are body-independent; the covariance is
$A^H\Cov(E(u))(A^H)^\top$.
More generally, for $s_{t+1}=A_{x_t}s_t+c_{x_t}$, if the symbol matrices
commute pairwise, equal symbol counts imply the same cue propagator
$\Phi_h=A_{x_{H-1}}\cdots A_{x_0}$ and the same centered cue covariance.
These are noiseless statements, also valid conditional on shared fixed noise.
\end{proposition}

The null concerns centered, body-conditioned cue information. It does not imply
equal native task error. Nor does commutativity eliminate noise effects:
intermediate noise is propagated through suffix products, which can differ
between equal-count bodies. With fixed $A$, cue-independent additive noise
with the same \emph{joint} law across arms produces the same terminal-noise law;
matching its one-step marginals alone is insufficient.

\begin{corollary}[Readout offsets can change error without changing memory]
\label{cor:offset}
Suppose two terminal states differ by a cue-independent translation and,
if noisy, have matched translated joint laws with the cue. A body-informed
decoder can compensate that translation and attain equal optimal risk.
For a fixed affine native head, let $e$ be its control-arm residual and
$\Delta$ its deterministic output shift. The squared-loss gap is
\begin{equation}
 \E\norm{e+\Delta}^2-\E\norm{e}^2
 =2\E[e]^\top\Delta+\norm{\Delta}^2.
 \label{eq:offset-gap}
\end{equation}
\end{corollary}

Thus an error contrast can identify readout sensitivity while leaving cue
recoverability unchanged. Conversely, noncommuting transitions permit, but do
not guarantee, unequal cue products or unequal recovery risk. The standard
adjacent-swap bound in \cref{lem:swaps} quantifies a possible product change;
it neither assigns a harmful direction to periodic order nor bounds decoder
error without further assumptions.

\subsection{Dimension, contraction, and recovery are distinct}

Consider $u\sim\mathcal N(0,I_d)$, stored as $s_H=a^Hu$ with $0<a<1$.
At every finite horizon, the signal covariance has participation ratio $d$,
although its trace is $d a^{2H}$. If the measured state is
$y_H=s_H+\varepsilon$ with independent
$\varepsilon\sim\mathcal N(0,\sigma^2I_d)$ and $\sigma>0$, its covariance
also has participation ratio $d$, while the optimal squared recovery risk is
\begin{equation}
 \E\norm{u-\E[u\mid y_H]}^2
 =\frac{d\sigma^2}{a^{2H}+\sigma^2}\longrightarrow d.
 \label{eq:attenuation-risk}
\end{equation}
Without that noise, $a^{-H}s_H$ recovers $u$ exactly. A scale-invariant
dimension statistic therefore cannot identify reliable memory loss; a noise
or precision model is essential. The converse is also possible:
$s_H=(u,a^Hv)$ for independent isotropic $u\in\R^d,v\in\R^k$ retains $u$
exactly while participation ratio decreases from $d+k$ toward $d$.

Local contraction is likewise compatible with finite categorical memory.
This is the elementary counterpart of phase-locked memory in driven recurrent
networks \citep{pals2024-phase-locked-memory}.
\begin{proposition}[Contracting categorical memory]
\label{prop:basins}
Let $\phi_t=0.2(-1)^t$, $r_t=s_t-\phi_t$, and
$s_{t+1}=g(r_t)+\phi_{t+1}$ with $g(r)=1.5r-0.5r^3$.
The two period-two responses $s_t=\pm1+\phi_t$ have zero tangent multipliers.
The tubes $|r_t\mp1|\le0.05$ are invariant and contracting, yet
$\operatorname{sign}(s_t)$ recovers their binary label for every $t$.
\end{proposition}

Contraction within a basin does not measure separation between basins. Exact
smooth memory of a continuum has a different local requirement, familiar from
continuous-attractor models \citep{seung1996}.
\begin{proposition}[Exact smooth memory requires neutral directions]
\label{prop:continuous}
Let $U\subset\R^d$ be open and let $E:U\to\R^n$ and
$\mathcal D:\R^n\to\R^d$ be $C^1$, with $\mathcal D(E(u))=u$.
Then $\rank DE(u)=d$. If a $C^1$ period map $\mathcal P$ also satisfies
$\mathcal P(E(u))=E(u)$ on $U$, then
$D\mathcal P(E(u))DE(u)=DE(u)$: its tangent map has at least $d$
independent eigenvectors with eigenvalue one.
\end{proposition}

This is an exact preserved-family statement. Finite-sample mean-squared
competence alone supplies neither a local singular-value bound nor a neutral
multiplier. Finally, a forced periodic response can persist with strictly
contracting multipliers (\cref{prop:scalar}), and stable eigenvalues allow
transient gain (\cref{prop:nonnormal}). Periodic activity, covariance dimension,
and tangent stability each require a separate task-memory measurement.
